% revised_thesis.tex — parallel build for side-by-side diff with thesis.tex
% Compile from project root: pdflatex revised_thesis && bibtex revised_thesis && pdflatex revised_thesis x2

\documentclass[12pt]{report}

\usepackage[language=english, type=bachelor]{style}
\usepackage{amsmath}
\usepackage{amssymb}

\newcommand{\thesisfig}[2]{%
  \IfFileExists{#2}{%
    \includegraphics[width=#1]{#2}%
  }{%
    \fbox{\parbox{#1}{\centering\vspace{1.1cm}\footnotesize\textit{Add file:} \texttt{\detokenize{#2}}\vspace{1.1cm}}}%
  }%
}

\usepackage{fancyhdr}
\setlength{\headheight}{27.35799pt}
\pagestyle{fancy}
\lhead{}
\chead{}
\renewcommand{\headrulewidth}{0.2pt}
\renewcommand{\footrulewidth}{0.2pt}

\begin{document}

\specialization{Computer Science}
\title{Augmented Reality-Based Interior Design: A Mobile Application for Parametric Room Scanning and Interactive 3D Scene Reconstruction}
\author{Mara Ielciu}
\supervisor{Professor Dr. Darabant Sergiu Adrian}

\maketitle

\newpage
\thispagestyle{empty}
\mbox{}
\newpage
\pagenumbering{roman}

\cleardoublepage
ABSTRACT
\vspace{0.5cm}
\hrule
\vspace{0.5cm}

This bachelor thesis presents an iOS application for augmented reality--based interior design framed as a unified \textbf{scan-to-stage engineering pipeline}. Users capture a real room with Apple's RoomPlan API to obtain a metric, semantically labelled \texttt{CapturedRoom}, then edit the space in a persistent, non-AR RealityKit editor: replacing detected furniture with USDZ models from a catalogue or from on-device Object Capture photogrammetry, styling walls and floors, and adding virtual decor and lighting. Room analytics derive measurements and renovation estimates from scan geometry; rule-based Smart Tips advise on lighting density and spatial balance. Completed designs export to Blender as glTF-based ZIP archives. For visual benchmarking, the app imports offline-trained 3D Gaussian Splatting \texttt{.ply} files and can request feedforward SHARP reconstructions from a user-controlled Mac on the local network, rendering both through a custom Metal splat viewer. A three-way comparison (mesh versus offline splats versus SHARP) clarifies trade-offs among geometric editability, capture effort, and plenoptic fidelity.

\tableofcontents

\newpage
\pagenumbering{arabic}

\section*{Context and Related Work}
\addcontentsline{toc}{section}{Context and Related Work}

Indoor scene understanding research, exemplified by ScanNet~\cite{dai2017scannet}, established that semantic three-dimensional structure enables spatial reasoning beyond raw geometry. Consumer AR furniture apps often lack persistent room models; K\'an~et~al.~\cite{kan2021ar} add layout intelligence but not off-site editing tied to a captured metric room. Kandil~et~al.~\cite{kandil2021ux} show that manipulation precision and session persistence are critical usability gaps in handheld AR design tools---gaps this project addresses with a non-AR editor and local JSON persistence.

\input{chapters/revised_chapter1_introduction}
\input{chapters/revised_chapter2_capture_theory}
\input{chapters/revised_chapter3_gaussian_splatting}
\input{chapters/revised_chapter4_requirements}
\input{chapters/revised_chapter5_implementation}
\input{chapters/revised_chapter6_conclusions}

% Optional: uncomment to include extended survey chapter
% \input{chapters/revised_state_of_the_art}

\bibliography{references}

\end{document}
